% CHAPTER 4
\chapter{Implementation}

\section{Hardware and Software Requirements}
[Specify all computational hardware and software environments used during development, training, and deployment:]
\begin{itemize}
    \item \textbf{Workstation Hardware:} [CPU processor model, GPU accelerator with VRAM, system RAM, storage type.]
    \item \textbf{Target / Edge Hardware:} [If applicable: embedded processor, mobile device, or edge SBC specs.]
    \item \textbf{Operating System:} [e.g., Ubuntu 22.04 LTS, Windows 11, macOS.]
    \item \textbf{Software Stack:} [Python version, CUDA/cuDNN versions, driver versions.]
\end{itemize}

\section{Tools and Frameworks}
Table \ref{tab:tools} outlines the primary libraries, development toolkits, and platforms utilized in this project:

\begin{table}[htbp]
\centering
\caption{Software frameworks and toolkits utilized in the implementation}
\label{tab:tools}
\vspace{0.2cm}
\resizebox{0.95\columnwidth}{!}{
\begin{tabular}{|l|l|l|}
\hline
\textbf{Framework / Library} & \textbf{Version} & \textbf{Functional Purpose in Project} \\
\hline
PyTorch / TensorFlow & [Version] & Core deep learning modeling, gradient computation, and backpropagation. \\
\hline
OpenCV / PIL & [Version] & Image preprocessing, transformation, and visual rendering utilities. \\
\hline
scikit-learn & [Version] & Dataset partitioning, baseline modeling, and evaluation metric computation. \\
\hline
Next.js / Streamlit / Flask & [Version] & Web interface / command dashboard for user interaction and demo. \\
\hline
\end{tabular}
}
\end{table}

\section{Dataset Partitioning}
[Describe how the dataset was split into training, validation, and testing subsets (e.g., 70:15:15 or 80:10:10). Detail whether the split was purely randomized, class-stratified, or subject/patient-wise to prevent data leakage.]

\section{Hyperparameter Settings}
Table \ref{tab:hyperparameters} summarizes the exact hyperparameter values configured during model training:

\begin{table}[htbp]
\centering
\caption{Hyperparameter settings configured for model training}
\label{tab:hyperparameters}
\vspace{0.2cm}
\begin{tabular}{|l|l|}
\hline
\textbf{Hyperparameter} & \textbf{Configured Value} \\
\hline
Input Resolution & [$H \times W \times C$, e.g., $640 \times 640 \times 3$] \\
\hline
Epochs & [e.g., 100] \\
\hline
Batch Size & [e.g., 16 or 32] \\
\hline
Optimizer & [e.g., AdamW / SGD] \\
\hline
Initial Learning Rate & [e.g., 0.001] \\
\hline
Learning Rate Schedule & [e.g., Cosine Annealing with Warmup] \\
\hline
Weight Decay / Regularization & [e.g., 0.0005] \\
\hline
\end{tabular}
\end{table}

\section{Model Training Procedure}
[Provide a step-by-step description of the training workflow: batch loading, forward pass, loss calculation, backpropagation, validation monitoring, and checkpoint saving. State the total training time and convergence epoch.]

\section{Deployment and User Interface}
[Describe the deployment target: web application, mobile app, edge device daemon, or API service. Detail how end-users or domain practitioners interact with the system, and include illustrative UI screenshots or system outputs.]
